\documentclass{article}

\usepackage[final]{neurips_2025}
\usepackage[utf8]{inputenc}
\usepackage[T1]{fontenc}
\usepackage{hyperref}
\usepackage{url}
\usepackage{booktabs}
\usepackage{amsfonts}
\usepackage{nicefrac}
\usepackage{microtype}
\usepackage{xcolor}
\usepackage{amsmath}
\usepackage{graphicx}

\title{Element Graph for Academic Figure Retrieval:\\
Graded Supervision and Encoder-agnostic Score Propagation\\
from a Single Edge-Weight Schema}

\author{%
  Jaehyeon Lee \\ 20262670
  \And
  Seoyeon Lee \\ 20262097
  \And
  Suhyeon Jun \\ 20262717
  \And
  Minjun Kim \\ 20262242
}

\begin{document}
\maketitle

\begin{abstract}
Academic papers contain rich cross-element evidence---a query is often answered
by a figure, its caption, the citing paragraph, and an appendix reference jointly,
linked through typed relations such as \texttt{caption\_of} and \texttt{refer\_to}.
We propose using the document's typed element graph as a \emph{single source of
truth} for both retrieval-time roles: (i)~a graph-distance kernel defines a
graded relevance distribution that supervises a listwise contrastive
loss (\textbf{GRCL}), strictly generalizing binary InfoNCE; (ii)~the same edge-weight
schema constructs a transition matrix for inference-time score propagation that
plugs into any base retriever. On SPIQA test-A, GRCL improves R@10 by $+2.7$pp
and MRR by $+5.4$pp over InfoNCE ($p\!=\!0.042$). Restricting GRCL to
\texttt{refer\_to} edges alone gives R@10 $=$ \textbf{90.4}, the best of our
43-row ablation. Most strikingly, applying the same graph propagation to the
external retrieval-tuned MLLM GME-Qwen2-VL-2B raises its SPIQA R@10 from $58.3$
to $\mathbf{84.4}$ ($+26$pp) and its MRR from $33.1$ to $\mathbf{63.3}$
($+30$pp)---the largest single effect in the study---without any fine-tuning.
Three propagation regimes (uniform diffusion, weighted diffusion, and Personalized
PageRank) all converge to R@10 $83$--$85\%$ on top of GME, indicating that
\emph{graph structure itself}, not modifier sophistication, is the source of
utility.
\end{abstract}

\section{Introduction}

The answer to a query on an academic paper is usually a \emph{set} of
related elements rather than a single one. For example,
``How does Method A improve over the baseline on the hard split, and how does
the appendix justify the gap?'' is answered by (i) Table~1, (ii) its caption
that defines ``hard split,'' (iii) a paragraph in Section~4.2 that references
the table, and (iv) an Appendix~B paragraph that the body cites. These
elements span multiple pages and modalities, yet are connected by
\emph{explicit typed edges}: \texttt{caption\_of}, \texttt{refer\_to}, and
body--appendix references. Retrieval evidence in academic documents is thus
intrinsically graph-structured.

Existing systems do not treat this graph as a first-class signal. Text-centric
RAG drops figures. Page-level visual RAG \cite{colpali,visrag} can miss
cross-page evidence. Element-level retrieval without graph priors ranks each
candidate independently and fails to recover references that bind elements
together.

We propose using the document's typed element graph as a \emph{single
source of truth} for retrieval. A single edge-weight schema drives two
independent contributions:

\textbf{(C1) Graph-Relevance Contrastive Loss (GRCL).} A 2-hop max-product
path kernel $g(q,e)$ defines a graded relevance distribution over candidates;
a listwise cross-entropy fits the retriever to it. GRCL strictly generalizes
binary InfoNCE (recovered when $r_i \in \{0,1\}$).

\textbf{(C2) Inference-time graph propagation as a plug-in.} The \emph{same}
edge weights construct a transition matrix $P$ used by weighted diffusion or
Personalized PageRank to smooth any base retriever's scores. No retraining
is required.

Our most striking empirical finding is that applying (C2) to the external
retrieval-tuned MLLM GME-Qwen2-VL-2B \cite{gme} raises its SPIQA R@10 from
$58.3$ to $84.4$ ($+26$pp) and almost doubles its MRR
($33.1\!\to\!63.3$)---the same schema simultaneously supervises a 200M
dual-encoder and augments a 2B decoder MLLM at inference time, evidence that
the graph functions as an \emph{encoder-agnostic universal prior}.

\section{Method}

\paragraph{Element graph schema.} Each document is a per-document graph
$G_d=(V_d,E_d)$ with six node types (\texttt{text}, \texttt{figure},
\texttt{table}, \texttt{caption}, \texttt{equation}, \texttt{section\_header})
and five edge types. The framework's \emph{single source of truth} is the
edge-weight schema:
\[
\text{BASE\_EDGE\_WEIGHTS} = \big\{\text{caption\_of}\!:\!1.0,\ \text{refer\_to}\!:\!0.8,\ \text{contains}\!:\!0.5,\ \text{reading\_next}\!:\!0.3,\ \text{section\_next}\!:\!0.2\big\}.
\]
A section-role modifier (appendix: $1.3$, references: $0.5$) and a
cross-modal visual-target modifier ($1.1$ for \texttt{refer\_to} into a
figure/table) refine the base weights based on the target node.

\paragraph{Encoder and late interaction.} We use SigLIPv2-base-patch16-224
\cite{siglip2} with LoRA rank~8 on attention projections (491k trainable
parameters, $0.32\%$ of the base). Token-level outputs are projected to
$D_\text{proj}\!=\!128$ and L2-normalized. The retrieval score uses
ColBERT-style MaxSim \cite{colbert}:
$\textstyle\mathrm{score}(q,e) = \sum_i \max_j \langle q_i, e_j\rangle$.

\paragraph{Graph-Relevance Contrastive Loss.} The graph relevance kernel
between an anchor $q$ and a candidate $e$ is the max-product path weight up
to two hops with hop decay $\gamma\!=\!0.5$:
\[
g(q,e) = \max_{\pi:q\to e,\,|\pi|\le 2}\ \prod_{(u,v)\in\pi} w_{uv}\cdot \gamma^{|\pi|-1},
\]
where $w_{uv}$ uses BASE\_EDGE\_WEIGHTS and target-modifier corrections.
For a candidate pool of $K\!=\!12$ elements (graph 1-hop positives plus
same-document and cross-document negatives), let $r_i\!=\!g(q,e_i)$ and
$s_i\!=\!\mathrm{score}(q,e_i)$. The GRCL listwise loss is
\[
\mathcal{L}_\text{GRCL}(q) = -\sum_i \frac{r_i}{\sum_j r_j}\ \log\frac{\exp(s_i/\tau)}{\sum_j \exp(s_j/\tau)}.
\]
When the targets degenerate to $\{0,1\}$, this recovers binary
InfoNCE---hence GRCL is a strict graded relaxation.

\paragraph{Inference-time graph propagation.} The base retriever's top-50
scores are propagated on the same per-document graph. The transition matrix
$P$ is the row-normalized weight matrix built from the \emph{identical}
schema above. We consider three regimes: (A) \emph{uniform diffusion} with all
edge weights set to one (structure-only), (B) \emph{weighted diffusion}
$s^{(t+1)} = (1\!-\!\alpha)s^{(t)} + \alpha P^\top s^{(t)}$, and
(C) \emph{Personalized PageRank} $s^{(t+1)} = \alpha q + (1\!-\!\alpha) P^\top s^{(t)}$
with teleport vector $q = s^{(0)}$ \cite{ppr}. All three are eval-only,
requiring no parameter updates.

\section{Experiments}

\paragraph{Dataset.} SPIQA \cite{spiqa} is a figure-grounded QA benchmark over
arXiv CS papers. We train on the train\_val split ($25{,}459$ papers,
$\sim$$250$k caption--figure pairs) and evaluate on test-A ($118$ papers,
$999$ queries). Graphs are extracted directly from arXiv LaTeX source
($\backslash\!$section, $\backslash\!$ref) and the dataset's
\texttt{all\_figures} JSON; no PDF parsing is required.

\paragraph{Training.} Adam\-W, $8{,}000$ steps, batch of $16$ anchors $\times$ $12$
candidates, $\gamma\!=\!0.5$. We compare two settings: default
$\text{lr}\!=\!3\mathrm{e}\text{-}5$, $\tau\!=\!0.07$; and a tuned
$\text{lr}\!=\!1\mathrm{e}\text{-}4$, $\tau\!=\!0.10$. Anchors are sampled
$1/3$ each from \texttt{caption\_of}, \texttt{refer\_to}, and natural-language
QA. One row trains in $\sim$$5$ wall-hours on an A100 80GB.

\paragraph{Evaluation.} Per-document candidate pool. We report R@5, R@10, MRR.
Statistical significance is assessed with paired bootstrap ($1{,}000$
resamples, $95\%$ CI).

\paragraph{Baseline.} (a) InfoNCE: the same encoder and training data but
GRCL is replaced by standard binary InfoNCE. (k=GME) GME-Qwen2-VL-2B-Instruct
zero-shot, an external 2B retrieval-tuned MLLM \cite{gme}.

\section{Results}

Our results are summarized in Table~\ref{tab:main}. We highlight five
findings, each supported by row references into the table.

\textbf{GRCL outperforms InfoNCE on in-domain SPIQA (rows 1 vs.\ 2).}
Graded graph supervision improves R@10 by $+2.7$pp and MRR by $+5.4$pp.
Paired bootstrap on Coverage@10 yields $\Delta\!=\!+0.0255$, $95\%$ CI
$[+0.0015,+0.0511]$, $p\!=\!0.042$. The relative MRR gain ($+16.5\%$) is
about $5\times$ the relative R@10 gain, indicating that graded supervision
sharpens \emph{top-rank quality} rather than merely recall coverage.

\textbf{The \texttt{refer\_to} edge dominates the graph signal (row 3).}
Restricting the GRCL kernel to a single edge type, \texttt{refer\_to}-only
yields R@10 $=$ \textbf{90.4}---the best of all $43$ ablation rows---a
$3.5$pp improvement over using all five edges equally. Explicit references
where a paragraph cites a figure carry the strongest training signal;
trivial adjacencies (\texttt{caption\_of}) and structural containment
(\texttt{contains}) are weaker.

\textbf{Hyperparameter tuning is a comparable lever (rows 4--5).}
Switching $\text{lr}$ from $3\mathrm{e}\text{-}5$ to $1\mathrm{e}\text{-}4$
alone gives $+1.7$pp R@10 / $+5.2$pp MRR. The full combo
($\text{lr}\!=\!1\mathrm{e}\text{-}4$, $\tau\!=\!0.10$,
$\lambda_\text{cov}\!=\!0.5$) reaches R@10 $88.0$ / MRR $42.2$.

\textbf{Graph propagation plugs into external MLLMs for the largest single
gain (rows 7 vs.\ 8--12).} Applying \emph{exactly the same} edge-weight
schema as inference-time score propagation on top of GME-Qwen2-VL-2B raises
its SPIQA R@10 from $58.3$ to $84.4$ ($+26$pp) and its MRR from $33.1$ to
$63.3$ ($+30$pp). Without any fine-tuning, GME$+$propagation reaches $96\%$
of our best trained SigLIP setting's R@10---demonstrating that the graph
schema is encoder-agnostic and that the propagation step is a deployable
plug-in.

\textbf{Three propagation regimes are equivalent on GME (rows 8--12).}
Uniform diffusion, weighted diffusion (the BASE $\times$ role $\times$ visual
modifier), and PPR all converge to R@10 $83$--$85\%$ and MRR $59$--$63\%$.
This indicates that the utility of propagation comes from the
\emph{existence} of the graph structure rather than from modifier
sophistication---a useful negative result for system designers.

\begin{table}[t]
\centering
\small
\caption{SPIQA test-A retrieval results. Bold cells mark the best in their
column or the headline finding. $\Delta$ values are versus the in-group baseline
(row~1 for trained encoders, row~7 for GME).}
\label{tab:main}
\begin{tabular}{lccccc}
\toprule
Setup & R@5 & R@10 & MRR & $\Delta$R@10 & $\Delta$MRR \\
\midrule
\multicolumn{6}{l}{\emph{Trained dual-encoder: SigLIPv2-base + LoRA, GRCL unless noted}} \\
1.\ (a) InfoNCE baseline & --- & 84.4 & 32.8 & --- & --- \\
2.\ (e) GRCL ($p\!=\!0.042$ vs.\ 1) & --- & 87.1 & 38.2 & $+2.7$ & $+5.4$ \\
3.\ + \texttt{refer\_to}-only edges & --- & \textbf{90.4} & 41.9 & $\mathbf{+6.0}$ & $+9.1$ \\
4.\ + $\text{lr}\!=\!1\mathrm{e}\text{-}4$ & --- & 88.6 & \textbf{42.6} & $+4.2$ & $\mathbf{+9.8}$ \\
5.\ + best HP combo (lr $\!+\!\tau\!+\!\lambda$) & --- & 88.0 & 42.2 & $+3.6$ & $+9.4$ \\
6.\ + CLIP-L/14 backbone swap & --- & 88.1 & 36.2 & $+3.7$ & $+3.4$ \\
\midrule
\multicolumn{6}{l}{\emph{External MLLM: GME-Qwen2-VL-2B-Instruct, zero-shot, no fine-tuning}} \\
7.\ GME, no propagation & 47.1 & 58.3 & 33.1 & --- & --- \\
8.\ \ + uniform diffusion ($\alpha\!=\!0.5,\,T\!=\!2$) & --- & 83.9 & 62.1 & $+25.6$ & $+29.0$ \\
9.\ \ + weighted diffusion ($\alpha\!=\!0.3,\,T\!=\!2$) & \textbf{78.2} & 84.4 & 63.3 & $+26.1$ & $+30.2$ \\
10.\ + weighted diffusion ($\alpha\!=\!0.3,\,T\!=\!3$) & --- & \textbf{84.5} & \textbf{63.4} & $\mathbf{+26.2}$ & $\mathbf{+30.3}$ \\
11.\ + PPR ($\alpha\!=\!0.30$ teleport) & --- & 83.3 & 59.1 & $+25.0$ & $+26.0$ \\
12.\ + PPR ($\alpha\!=\!0.50$) & --- & 83.2 & 58.6 & $+24.9$ & $+25.5$ \\
\bottomrule
\end{tabular}
\end{table}

\section{Conclusion}

A typed document element graph can serve simultaneously as graded
supervision for retriever training and as an inference-time prior for any
base retriever---both driven by a single edge-weight schema. On SPIQA,
GRCL gives a statistically significant R@10 $/$ MRR gain over binary
InfoNCE; restricting the graph kernel to \texttt{refer\_to} edges alone
delivers the best in-domain trained setting; and the schema's largest impact
is observed not in our own model but as a zero-shot plug-in atop the
external GME-Qwen2-VL-2B MLLM, where R@10 rises by $+26$pp and MRR by
$+30$pp. The equivalence of three propagation regimes (uniform diffusion,
weighted diffusion, Personalized PageRank) indicates that the source of
utility is the \emph{graph structure} rather than the modifier design.
We see two natural extensions: extending propagation to other
retrieval-tuned MLLMs as a general inference-time module, and exposing
edge-type weights as learnable parameters so that the dominance of
\texttt{refer\_to} can be discovered, not specified.

\bibliographystyle{plain}
\begin{thebibliography}{9}

\bibitem{colpali}
Faysse, M., Sibille, H., Wu, T., Omrani, B., Viaud, G., Hudelot, C., \& Colombo, P.
\textit{ColPali: Efficient Document Retrieval with Vision Language Models}. 2024.

\bibitem{visrag}
Yu, S., Tang, C., Xu, B., \textit{et al.}
\textit{VisRAG: Vision-based Retrieval-augmented Generation on Multi-modality Documents}. 2024.

\bibitem{gme}
Zhang, X., \textit{et al.}
\textit{GME: Improving Universal Multimodal Retrieval by Multimodal LLMs}. 2024.
\href{https://huggingface.co/Alibaba-NLP/gme-Qwen2-VL-2B-Instruct}{HF: Alibaba-NLP/gme-Qwen2-VL-2B-Instruct}.

\bibitem{siglip2}
Tschannen, M., Gritsenko, A., Wang, X., \textit{et al.}
\textit{SigLIP 2: Multilingual Vision-Language Encoders with Improved Semantic Understanding, Localization, and Dense Features}. 2025.

\bibitem{colbert}
Khattab, O.\ \& Zaharia, M.
\textit{ColBERT: Efficient and Effective Passage Search via Contextualized Late Interaction over BERT}. SIGIR 2020.

\bibitem{spiqa}
Pramanick, S., \textit{et al.}
\textit{SPIQA: A Dataset for Multimodal Question Answering on Scientific Papers}. 2024.

\bibitem{ppr}
Haveliwala, T.\ H.
\textit{Topic-Sensitive PageRank}. WWW 2002.

\bibitem{labelprop}
Zhou, D., Bousquet, O., Lal, T.~N., Weston, J., \& Schölkopf, B.
\textit{Learning with Local and Global Consistency}. NeurIPS 2003.

\end{thebibliography}

\end{document}
